% GENERATED FILE. Edit canonical lessons, then run npm run generate:books.
% canonical-source-hash: fnv1a64:35c278b5913d56e7
% canonical-lessons: SA-C394-tirthayatra, SA-C394-pratisthate, SA-C394-patheyam, SA-C394-yati, SA-C394-kamandaluh

\chapter{Setting Out on Pilgrimage}
\label{ch:sa-setting-out-on-pilgrimage}
\hlchaptermodality{394}

\begin{chapteropening}
\textbf{By the end of this chapter:} \emph{I can say when and how I set out on a pilgrimage.}

Complete the last lesson of chapter 394: I can say when and how I set out on a pilgrimage.
\end{chapteropening}

\section[\texorpdfstring{\sk{तीर्थयात्रा} (tīrthayātrā)}{तीर्थयात्रा (tīrthayātrā)}]{\sk{तीर्थयात्रा} (tīrthayātrā) — a pilgrimage}

\label{lesson:SA-C394-tirthayatra}



\begin{quote}
Before the new one: say the Sanskrit for radio, then the Sanskrit for the internet.
\end{quote}

\subsection*{You'll want to know: \sk{तीर्थयात्रा}}
\textbf{\sk{तीर्थयात्रा}} — \emph{tīrthayātrā} — \textquotedblleft{}a pilgrimage\textquotedblright{}.

\begin{quote}
\sk{अहं} \sk{श्वः} \sk{तीर्थयात्रां} \sk{गमिष्यामि।}
\end{quote}

\begin{quote}
\emph{ahaṃ śvaḥ tīrthayātrāṃ gamiṣyāmi.}
\end{quote}

\begin{quote}
Tomorrow I will go on a pilgrimage.
\end{quote}

A feminine noun.

\begin{cousinweb}
\textbf{\sk{तीर्थयात्रा}} (\emph{tīrthayātrā}) is built from \emph{tīrtha}, \textquotedblleft{}a ford, a holy bathing place\textquotedblright{}, and \emph{yātrā}, \textquotedblleft{}a journey\textquotedblright{}; MW gives the compound as a journey to holy places, a pilgrimage (MW).
\end{cousinweb}

\begin{grammarlens}[title={What you've built}]
One more word for getting around.
\end{grammarlens}

\subsection*{Your turn}
\begin{itemize}
\raggedright
  \item \emph{Say it:} \emph{tīrthayātrā}
  \item \emph{Say it:} \emph{tīrthayātrā}, once more
  \item \emph{Say it:} \emph{ahaṃ śvaḥ tīrthayātrāṃ gamiṣyāmi.} — Tomorrow I will go on a pilgrimage.
  \item \emph{Recall:} say the Sanskrit for radio, then the Sanskrit for the internet, then say \emph{tīrthayātrā} again
\end{itemize}

\begin{tcolorbox}[breakable,colback=teal!4,colframe=teal!35!black,title={Before you move on}]
What does \sk{तीर्थयात्रा} mean? (\textquotedblleft{}A pilgrimage\textquotedblright{}.) Say it once more.
\end{tcolorbox}
\section[\texorpdfstring{\sk{प्रतिष्ठते} (pratiṣṭhate)}{प्रतिष्ठते (pratiṣṭhate)}]{\sk{प्रतिष्ठते} (pratiṣṭhate) — starts on a journey; sets out}

\label{lesson:SA-C394-pratisthate}



\begin{quote}
Before the new one: say the Sanskrit for the internet, then the Sanskrit for a pilgrimage.
\end{quote}

\subsection*{You'll want to know: \sk{प्रतिष्ठते}}
\textbf{\sk{प्रतिष्ठते}} — \emph{pratiṣṭhate} — \textquotedblleft{}starts on a journey; sets out\textquotedblright{}.

\begin{quote}
\sk{यात्रिकः} \sk{प्रातः} \sk{गृहात्} \sk{प्रतिष्ठते।}
\end{quote}

\begin{quote}
\emph{yātrikaḥ prātaḥ gṛhāt pratiṣṭhate.}
\end{quote}

\begin{quote}
The pilgrim sets out from home in the morning.
\end{quote}

Past (he …): \sk{सः} \sk{प्रातः} \sk{प्रस्थितवान्}, \emph{saḥ prātaḥ prasthitavān}.

\begin{cousinweb}
\textbf{\sk{प्रतिष्ठते}} (\emph{pratiṣṭhate}) is built from the prefix \emph{pra}, \textquotedblleft{}forth\textquotedblright{}, and the root √sthā, \textquotedblleft{}to stand\textquotedblright{}, which also gives \emph{tiṣṭhati}, \textquotedblleft{}stands\textquotedblright{}; MW gives \emph{pra-sthā} as to set out or depart, with the past participle \emph{prasthita}, and Pāṇini (1.3.22) gives √sthā the middle endings after \emph{pra} (MW).
\end{cousinweb}

\begin{grammarlens}[title={What you've built}]
One more word for getting around.
\end{grammarlens}

\subsection*{Your turn}
\begin{itemize}
\raggedright
  \item \emph{Say it:} \emph{pratiṣṭhate}
  \item \emph{Say it:} \emph{pratiṣṭhate}, once more
  \item \emph{Say it:} \emph{yātrikaḥ prātaḥ gṛhāt pratiṣṭhate.} — The pilgrim sets out from home in the morning.
  \item \emph{Recall:} say the Sanskrit for the internet, then the Sanskrit for a pilgrimage, then say \emph{pratiṣṭhate} again
\end{itemize}

\begin{tcolorbox}[breakable,colback=teal!4,colframe=teal!35!black,title={Before you move on}]
What does \sk{प्रतिष्ठते} mean? (\textquotedblleft{}Starts on a journey; sets out\textquotedblright{}.) Say it once more.
\end{tcolorbox}
\section[\texorpdfstring{\sk{पाथेयम्} (pātheyam)}{पाथेयम् (pātheyam)}]{\sk{पाथेयम्} (pātheyam) — food for the journey}

\label{lesson:SA-C394-patheyam}



\begin{quote}
Before the new one: say the Sanskrit for a pilgrimage, then the Sanskrit for starts on a journey.
\end{quote}

\subsection*{You'll want to know: \sk{पाथेयम्}}
\textbf{\sk{पाथेयम्}} — \emph{pātheyam} — \textquotedblleft{}food for the journey\textquotedblright{}.

\begin{quote}
\sk{माता} \sk{मह्यं} \sk{पाथेयं} \sk{ददाति।}
\end{quote}

\begin{quote}
\emph{mātā mahyaṃ pātheyaṃ dadāti.}
\end{quote}

\begin{quote}
Mother gives me food for the journey.
\end{quote}

A neuter noun.

\begin{cousinweb}
\textbf{\sk{पाथेयम्}} (\emph{pātheyam}) is built on \emph{pathin}, \textquotedblleft{}a road\textquotedblright{}, with the ending \emph{-eya}, \textquotedblleft{}belonging to\textquotedblright{}; MW gives \emph{pātheya} as provisions for a journey (MW).
\end{cousinweb}

\begin{grammarlens}[title={What you've built}]
One more word for getting around.
\end{grammarlens}

\subsection*{Your turn}
\begin{itemize}
\raggedright
  \item \emph{Say it:} \emph{pātheyam}
  \item \emph{Say it:} \emph{pātheyam}, once more
  \item \emph{Say it:} \emph{mātā mahyaṃ pātheyaṃ dadāti.} — Mother gives me food for the journey.
  \item \emph{Recall:} say the Sanskrit for a pilgrimage, then the Sanskrit for starts on a journey, then say \emph{pātheyam} again
\end{itemize}

\begin{tcolorbox}[breakable,colback=teal!4,colframe=teal!35!black,title={Before you move on}]
What does \sk{पाथेयम्} mean? (\textquotedblleft{}Food for the journey\textquotedblright{}.) Say it once more.
\end{tcolorbox}
\section[\texorpdfstring{\sk{याति} (yāti)}{याति (yāti)}]{\sk{याति} (yāti) — travels; goes}

\label{lesson:SA-C394-yati}



\begin{quote}
Before the new one: say the Sanskrit for starts on a journey, then the Sanskrit for food for the journey.
\end{quote}

\subsection*{You'll want to know: \sk{याति}}
\textbf{\sk{याति}} — \emph{yāti} — \textquotedblleft{}travels; goes\textquotedblright{}.

\begin{quote}
\sk{सः} \sk{धूमयानेन} \sk{नगरं} \sk{याति।}
\end{quote}

\begin{quote}
\emph{saḥ dhūmayānena nagaraṃ yāti.}
\end{quote}

\begin{quote}
He travels to the city by train.
\end{quote}

Past (he …): \sk{सः} \sk{नगरं} \sk{यातवान्}, \emph{saḥ nagaraṃ yātavān}.

\begin{cousinweb}
\textbf{\sk{याति}} (\emph{yāti}) comes from the root √yā, \textquotedblleft{}to go, to travel\textquotedblright{}, which also gives \emph{yātrā}, \textquotedblleft{}a journey\textquotedblright{}, and \emph{yāna}, \textquotedblleft{}a vehicle\textquotedblright{}; MW gives the participle \emph{yāta}, \textquotedblleft{}gone\textquotedblright{}; it lives on as Hindi \emph{jānā}, \textquotedblleft{}to go\textquotedblright{} (MW; CDIAL 10452).
\end{cousinweb}

\begin{grammarlens}[title={What you've built}]
One more word for getting around.
\end{grammarlens}

\subsection*{Your turn}
\begin{itemize}
\raggedright
  \item \emph{Say it:} \emph{yāti}
  \item \emph{Say it:} \emph{yāti}, once more
  \item \emph{Say it:} \emph{saḥ dhūmayānena nagaraṃ yāti.} — He travels to the city by train.
  \item \emph{Recall:} say the Sanskrit for starts on a journey, then the Sanskrit for food for the journey, then say \emph{yāti} again
\end{itemize}

\begin{tcolorbox}[breakable,colback=teal!4,colframe=teal!35!black,title={Before you move on}]
What does \sk{याति} mean? (\textquotedblleft{}Travels; goes\textquotedblright{}.) Say it once more.
\end{tcolorbox}
\section[\texorpdfstring{\sk{कमण्डलुः} (kamaṇḍaluḥ)}{कमण्डलुः (kamaṇḍaluḥ)}]{\sk{कमण्डलुः} (kamaṇḍaluḥ) — an ascetic's water-pot}

\label{lesson:SA-C394-kamandaluh}



\begin{quote}
Before the new one: say the Sanskrit for food for the journey, then the Sanskrit for travels.
\end{quote}

\subsection*{You'll want to know: \sk{कमण्डलुः}}
\textbf{\sk{कमण्डलुः}} — \emph{kamaṇḍaluḥ} — \textquotedblleft{}an ascetic's water-pot\textquotedblright{}.

\begin{quote}
\sk{यात्रिकः} \sk{कमण्डलौ} \sk{जलं} \sk{वहति।}
\end{quote}

\begin{quote}
\emph{yātrikaḥ kamaṇḍalau jalaṃ vahati.}
\end{quote}

\begin{quote}
The pilgrim carries water in his water-pot.
\end{quote}

A masculine noun.

\begin{cousinweb}
\textbf{\sk{कमण्डलुः}} (\emph{kamaṇḍaluḥ}) is the word MW gives for the water-pot that ascetics and pilgrims carry; MW gives no derivation, so its origin is uncertain (MW).
\end{cousinweb}

\begin{grammarlens}[title={What you've built}]
One more word for getting around.
\end{grammarlens}

\subsection*{Your turn}
\begin{itemize}
\raggedright
  \item \emph{Say it:} \emph{kamaṇḍaluḥ}
  \item \emph{Say it:} \emph{kamaṇḍaluḥ}, once more
  \item \emph{Say it:} \emph{yātrikaḥ kamaṇḍalau jalaṃ vahati.} — The pilgrim carries water in his water-pot.
  \item \emph{Recall:} say the Sanskrit for food for the journey, then the Sanskrit for travels, then say \emph{kamaṇḍaluḥ} again
\end{itemize}

\begin{tcolorbox}[breakable,colback=teal!4,colframe=teal!35!black,title={Before you move on}]
What does \sk{कमण्डलुः} mean? (\textquotedblleft{}An ascetic's water-pot\textquotedblright{}.) Say it once more.
\end{tcolorbox}
